% =====================================================
% 题型：外心向量代数余弦值选择题 (q_lyb_01_8_circumcenter_cos.tex)
% =====================================================
% =====================================================
% 难度：★★ (中档)
% =====================================================
\newcommand{\qVectorCircumcenterCosStem}{已知点 \(O\) 在 \(\triangle ABC\) 所在平面内，满足 \(|\vec{OA}| = |\vec{OB}| = |\vec{OC}|\)，且 \(\vec{AO} = 2\vec{AB} + 3\vec{AC}\)，则 \(\cos 2\angle BAC =\) \hfill （\quad）}

\newcommand{\qVectorCircumcenterCosOptions}{%
  \mcoptions[4]{\(-\dfrac{1}{4}\)}{\(-\dfrac{5}{8}\)}{\(\dfrac{3}{4}\)}{\(\dfrac{1}{4}\)}%
}

\newcommand{\qVectorCircumcenterCosAnswer}{D}

\newcommand{\qVectorCircumcenterCosSolution}{%
  设点 \(O\) 为原点，由 \(|\vec{OA}| = |\vec{OB}| = |\vec{OC}| = R\) 知，点 \(O\) 为 \(\triangle ABC\) 的外心，
  并且各顶点到 \(O\) 的距离等于外接圆半径 \(R\)。
  已知 \(\vec{AO} = 2\vec{AB} + 3\vec{AC}\)，将其用起点为 \(O\) 的向量表示：
  \(-\vec{OA} = 2(\vec{OB} - \vec{OA}) + 3(\vec{OC} - \vec{OA}) \implies 4\vec{OA} = 2\vec{OB} + 3\vec{OC}\)。
  对等式两边进行平方运算：
  \(|4\vec{OA}|^2 = |2\vec{OB} + 3\vec{OC}|^2 \implies 16|\vec{OA}|^2 = 4|\vec{OB}|^2 + 9|\vec{OC}|^2 + 12\vec{OB} \cdot \vec{OC}\)。
  将模长代入，得：
  \(16R^2 = 13R^2 + 12R^2\cos\angle BOC \implies 3R^2 = 12R^2\cos\angle BOC\)。
  由于 \(R > 0\)，消去 \(R^2\) 解得 \(\cos\angle BOC = \dfrac{1}{4}\)。
  由圆周角定理及外心的几何特征，圆心角 \(\angle BOC\) 与圆周角 \(\angle BAC\) 的数量关系满足：
  \(\angle BOC = 2\angle BAC\) （外心在三角形内部/同侧）
  或 \(\angle BOC = 2\pi - 2\angle BAC\) （外心在三角形外部/异侧）。
  在这两种情况下，恒有 \(\cos 2\angle BAC = \cos\angle BOC = \dfrac{1}{4}\)。
  故选 D。%
}


% =====================================================
% 别名兼容：旧课次宏定义
% =====================================================
\newcommand{\qLYBZeroOneEightStem}{\qVectorCircumcenterCosStem}
\newcommand{\qLYBZeroOneEightOptions}{\qVectorCircumcenterCosOptions}
\newcommand{\qLYBZeroOneEightAnswer}{\qVectorCircumcenterCosAnswer}
\newcommand{\qLYBZeroOneEightSolution}{\qVectorCircumcenterCosSolution}
